\bookchapter{บทที่ 64: การจัดเก็บข้อมูลและการเรียกใช้เครือข่าย}

\chapterintro{แอปที่ทำได้แค่ให้'ดู'เป็นแอปไร้ชีวิต แอปที่เก็บข้อมูลได้และต่อเน็ตได้ต่างหากที่มีชีวิตจริง บทนี้พาไปรู้จัก'สามพี่น้อง'ของชั้นข้อมูลใน iOS: UserDefaults (เก็บการตั้งค่าเล็ก ๆ), SwiftData (เก็บข้อมูลมีโครงสร้าง), URLSession (คุยกับเซิร์ฟเวอร์) พร้อมด้วย Codable สำหรับแปลง JSON และ async/await สไตล์ยุคใหม่ เรียนจบบทนี้ คุณจะเขียนแอปจริงที่จดจำนิสัยผู้ใช้และอัปเดตเนื้อหาแบบเรียลไทม์ได้แล้ว}

\booksection{เลือกวิธีจัดเก็บอย่างไร: ดูตารางนี้ก่อน}

\begin{center}
\begin{tabularx}{\linewidth}{|X|X|X|X|}
\hline
\textbf{ทางเลือก} & \textbf{เหมาะกับเก็บอะไร} & \textbf{จุดเด่น} & \textbf{ใช้ตรงไหนในหนังสือเล่มนี้} \tabularnewline
\hline
UserDefaults & การตั้งค่าเรื่อยเปื่อย (สวิตช์, สิ่งที่ชอบ) & คู่คีย์-ค่า ง่ายและตรงไปตรงมา & การตั้งค่าบทที่ 65, เมืองโปรดบทที่ 66 \tabularnewline
SwiftData & ข้อมูลมีโครงสร้าง (รายการบัญชี, บันทึกต่าง ๆ) & ตัวห่อฐานข้อมูลทางการ ประกาศด้วย @Model บรรทัดเดียว & บันทึกรายการบัญชีบทที่ 65 \tabularnewline
URLSession & (ไม่เก็บ ดึงผ่านเน็ต) ข้อมูลจากเซิร์ฟเวอร์ & เฟรมเวิร์กเครือข่ายมากับระบบ ไม่ต้องพึ่งไลบรารีภายนอก & ข้อมูลอากาศบทที่ 66 \tabularnewline
\hline
\end{tabularx}
\end{center}

\begin{center}\small\color{gray}หน้าที่ของสามพี่น้องฝ่ายข้อมูล\end{center}

จดคำกลับไว้ว่า'การตั้งค่าเล็ก ๆ เข้า UserDefaults ข้อมูลเป็นชุดใหญ่เข้า SwiftData ข้อมูลจากภายนอกใช้ URLSession' Core Data ในบทเรียนเก่ายังใช้ได้อยู่ (SwiftData ก็คือตัวตนยุคใหม่ของมัน) ส่วนการเขียน SQLite เองแทบไม่ต้องแตะ — เรียนของใหม่ตรง ๆ ดีกว่า อย่าเดินซ้ำทางเก่า

\booksection{UserDefaults: เก็บการตั้งค่าเล็ก ๆ น้อย ๆ}

\begin{lstlisting}[language=Swift]
// เขียน: เก็บค่า
UserDefaults.standard.set(true, forKey: "isDarkMode")
UserDefaults.standard.set("ปักกิ่ง", forKey: "favoriteCity")

// อ่าน: ดึงค่าออกมา (ไม่มีก็คืนค่าเริ่มต้น)
let isDark = UserDefaults.standard.bool(forKey: "isDarkMode")
let city = UserDefaults.standard.string(forKey: "favoriteCity") ?? "ปักกิ่ง"

// ลบ: เอาคีย์นี้ออก
UserDefaults.standard.removeObject(forKey: "isDarkMode")
\end{lstlisting}

ระวังว่ามันเหมาะกับการตั้งค่าเล็ก ๆ เท่านั้น — ถ้ามีรายการเป็นสิบเป็นร้อย ให้ส่งต่อให้ SwiftData จัดการ อย่าเอา UserDefaults มาใช้แทนฐานข้อมูล (จะทำให้แอปเปิดช้าลงและค้นหาไม่สะดวก)

\booksection{SwiftData: สร้าง'บ้านจริง ๆ'ให้ข้อมูล}

SwiftData คือแนวทางฐานข้อมูลทางการตั้งแต่ iOS 17 เป็นต้นมา ย่อ'สร้างตาราง เพิ่ม ลบ แก้ ค้น'ให้ง่ายสุดขีด — แค่ใส่แอนโนเทชัน @Model คลาสข้อมูลก็กลายเป็นตารางในฐานข้อมูล:

\begin{lstlisting}[language=Swift]
import SwiftData

// 1. @Model: คลาสนี้กลายเป็น "ตาราง" อัตโนมัติ
@Model
class Record {
    var amount: Double
    var category: String
    var note: String
    var date: Date

    init(amount: Double, category: String, note: String, date: Date) {
        self.amount = amount
        self.category = category
        self.note = note
        self.date = date
    }
}

// 2. ผูกคอนเทนเนอร์ที่จุดเริ่มแอป
@main
struct MyApp: App {
    var body: some Scene {
        WindowGroup {
            ContentView()
        }
        .modelContainer(for: Record.self)   // ฐานข้อมูลพร้อม
    }
}

// 3. ในวิวใช้ @Query อ่าน, @Environment เขียน
struct RecordListView: View {
    @Query(sort: \Record.date, order: .reverse) var records: [Record]
    @Environment(\.modelContext) private var context

    var body: some View {
        List(records) { r in
            Text("\(r.category) \(r.amount) หยวน")
        }
    }

    func add(amount: Double, category: String) {
        let r = Record(amount: amount, category: category, note: "", date: Date())
        context.insert(r)     // บันทึก
        try? context.save()   // ลงดิสก์
    }
}
\end{lstlisting}

เทมเพลต '@Model + modelContainer + @Query' ชุดนี้เอาไปวางในโปรเจกต์ไหนของคุณก็ได้ตรง ๆ — เปลี่ยนชื่อฟิลด์เท่านั้น พอข้อมูลเปลี่ยน @Query จะรีเฟรชอินเทอร์เฟซให้เอง ช่างมหัศจรรย์นัก

\booksection{Codable: คอย'แปลภาษา'ให้ JSON}

ข้อมูลจากเซิร์ฟเวอร์ส่วนใหญ่เป็นข้อความ JSON โปรโตคอล Codable ของ Swift ทำให้การแปลง'JSON struct'เกิดขึ้นเองอัตโนมัติ — ชื่อฟิลด์ตรงกัน มนต์ก็ทำงาน:

\begin{lstlisting}[language=Swift]
// JSON: {"name":"ปักกิ่ง","main":{"temp":26.5},"weather":[{"description":"แจ่มใส"}]}
struct WeatherResponse: Codable {
    let name: String
    let main: Main
    let weather: [Weather]
}
struct Main: Codable { let temp: Double }
struct Weather: Codable { let description: String }

// จากข้อความ JSON  struct
let decoder = JSONDecoder()
let data = jsonString.data(using: .utf8)!
let result = try decoder.decode(WeatherResponse.self, from: data)
print(result.weather[0].description)   // "แจ่มใส"
\end{lstlisting}

กฎเหล็ก: ดู JSON ที่ส่งกลับมาจริงก่อน แล้วค่อยนิยามฟิลด์ บั๊ก'แปลงค่าไม่สำเร็จ'ส่วนใหญ่มาจากชื่อฟิลด์พิมพ์ผิดหรือเดาชนิดข้อมูลผิด (เซิร์ฟเวอร์ส่ง “26.5” เป็นสตริง แต่คุณประกาศเป็น Double) — พรินต์ JSON ดิบออกมาเทียบดู หาจุดผิดเจอในสิบวินาที

\booksection{URLSession + async/await: คุยกับเซิร์ฟเวอร์}

การต่อเน็ตบน iOS ใช้ URLSession ที่มากับระบบ คู่กับ async/await (ฟังก์ชันอะซิงโครนัส) ของ Swift ทำให้'การรอเน็ต'เป็นธรรมชาติเหมือน'รอพัสดุ':

\begin{lstlisting}[language=Swift]
struct WeatherService {
    // ฟังก์ชัน async: async หมายถึง "ต้องรอสักครู่" throws หมายถึง "อาจล้มเหลว"
    func fetchWeather(city: String) async throws -> WeatherResponse {
        let url = URL(string:
            "https://api.example.com/weather?city=\(city)")!
        let (data, _) = try await URLSession.shared.data(from: url)
        return try JSONDecoder().decode(WeatherResponse.self, from: data)
    }
}

// จุดเรียก (ใน .task):
.task {
    do {
        let weather = try await service.fetchWeather(city: "ปักกิ่ง")
        // อัปเดตหน้าจอ……
    } catch {
        print("เกิดข้อผิดพลาด: \(error)")
    }
}
\end{lstlisting}

\textbf{กันพลาด: กฎเหล็กสองข้อของการใช้เครือข่าย}\par

(1) การต่อเน็ตต้องใช้ async/await ห้ามปิดกั้นเธรดหลักเด็ดขาด — อินเทอร์เฟซค้างคือประสบการณ์ระดับทำให้แอปล้ม โครงสร้างคอนเคอเรนซีของ Swift ทำให้เรื่องนี้เป็นธรรมชาติมาก แต่อย่าลืม try await กับ do/catch (2) ATS ห้าม HTTP แบบไม่เข้ารหัสเป็นค่าเริ่มต้น: iOS บังคับให้ใช้ HTTPS ตอนทดสอบบริการ http:// ในเครื่องต้องเปิดข้อยกเว้นชั่วคราวใน Info.plist (NSAppTransportSecurity) แต่ผลิตภัณฑ์จริงต้องใช้ HTTPS ทั้งหมด — ทีมรีวิวก็เฝ้าดูเรื่องนี้เช่นกัน

\begin{tcbpractice}{ลงมือทำ: สร้างแอปเล็ก ๆ 'คำคมประจำวัน'}
\booksection{ข้อผิดพลาดเครือข่ายก็ต้อง'จัดระดับ': อย่าโยน error ทุกชนิดให้ผู้ใช้}

เน็ตหลุด, หมดเวลา, 404, เซิร์ฟเวอร์ 500... แต่ละความผิดพลาด ผู้ใช้ควรเห็นข้อความที่ไม่เหมือนกัน:

\begin{lstlisting}[language=Swift]
func friendlyMessage(_ error: Error) -> String {
    if let urlError = error as? URLError {
        switch urlError.code {
        case .notConnectedToInternet: return "เน็ตหลุด ลองเช็คการเชื่อมต่อไหม?"
        case .timedOut:               return "คำขอหมดเวลา ลองใหม่ภายหลังไหม?"
        case .cannotFindHost:         return "หาเซิร์ฟเวอร์ไม่เจอ"
        default:                      return "เน็ตมีปัญหา: \(urlError.code)"
        }
    }
    if let d = error as? DecodingError {
        return "รูปแบบข้อมูลไม่ตรง (อาจเพราะ API เปลี่ยน)"
    }
    return "เกิดข้อผิดพลาด: \(error.localizedDescription)"
}
\end{lstlisting}

แสดง'ภาษาคน'ให้ผู้ใช้ ส่วน'ต้นฉบับ'ปล่อยไว้ในล็อก (พรินต์ error ฉบับเต็ม) อย่าเอาสำนวนอังกฤษกึกกือจาก error.localizedDescription โยนใส่ผู้ใช้ตรง ๆ — ความเป็นแอปชั้นดีซ่อนอยู่ในรายละเอียดพวกนี้

หา API คำคมฟรีสักตัว (อย่าง hitokoto หรือแนวเดียวกัน) ใช้ URLSession + async/await ดึง'คำคมประจำวัน'มาแสดงกลางหน้าจอ เพิ่มปุ่ม'เปลี่ยนประโยค' แล้วเก็บประโยคที่ชอบลง'รายการโปรด'ด้วย SwiftData แบบฝึกหัดนี้ใช้ครบทั้ง UserDefaults/SwiftData/Codable/URLSession — ทำเสร็จแล้ว บทที่ 7 กับ 8 ก็ง่ายเหมือนกินขนม

\end{tcbpractice}

\begin{tcbsummary}{สรุปบทนี้}
\begin{itemize}
  \item คำกลับ: การตั้งค่าเล็ก ๆ เข้า UserDefaults ข้อมูลเป็นชุดใหญ่เข้า SwiftData ข้อมูลจากภายนอกใช้ URLSession
  \item สามท่าเด็ดของ SwiftData: @Model ประกาศตาราง, modelContainer ติดตั้งคอนเทนเนอร์, @Query สำหรับอ่าน + context.insert สำหรับเขียน
  \item Codable'แปลภาษา'JSON ให้เองอัตโนมัติ: ชื่อฟิลด์ตรงกัน มนต์ก็ทำงาน
  \item URLSession + async/await: ต่อเน็ตเหมือนรอพัสดุ; do/catch จัดการความล้มเหลว
  \item กฎเหล็กสองข้อ: อย่าปิดกั้นเธรดหลัก; ผลิตภัณฑ์จริงต้องใช้ HTTPS (ATS)
\end{itemize}

\end{tcbsummary}
